\section{Discussion}
Familiarity-dependent inhibition improved exploration before the motor readout learned from reward. Both calcium memory and the exponential trace increased coverage and reduced immediate reversals. After reinforcement learning, memory was most beneficial on Obstacles and Branching, whereas the learned controllers approached the reward ceiling on Open without it. Broader exploration did not always improve goal-directed navigation: on Open, Trace explored more of the map but obtained a lower mean return than the uniform controls.

The calcium module provides a transient record of spatial experience that can be queried through predictions of candidate actions. Its competitive performance with the neuronal trace offers computational support for the hypothesis that astrocytic dynamics could fulfill trace-like memory functions. Trace had higher mean coverage without RL on every map, while the ordering of the memory variants after RL depended on the environment. Because each controller was tuned independently, differences in thresholds, motor calibration, and optimization also contribute to this comparison. The experiments support the usefulness of persistent familiarity memory without identifying a consistent calcium-specific advantage.

Removing memory from trained controllers reduced returns on the structured maps, particularly Obstacles and Branching. Memory thus remained part of action selection after reward learning; its contribution was not confined to guiding exploration during training. Disabling memory does not reproduce training without memory, and the resulting performance loss does not identify how the actor has adapted to inhibition. Whether the actor learns stronger excitation to compensate for inhibition remains an open question.

The detached architecture separates the signals that modify the recurrent network and motor readout. Prediction error trains the recurrent representation, while reward trains the actor. Their interaction continues through behavior: policy changes alter the transitions available for prediction learning, and continued prediction learning changes the representation received by the actor. Detachment therefore prevents direct reward-driven changes to recurrent synapses without making the learning processes independent. A matched non-detached control is needed to determine whether this separation improves performance.

Prediction on a controller's own trajectories should also be distinguished from prediction across shared test contexts. High next-cell accuracy along successful routes can coexist with errors on less frequently visited transitions. The common-input probe exposes this difference, although its weighting of moving and collision transitions also differs from the trajectory-based measure. Moreover, the controller queries memory using predicted spatial features, whereas next-cell accuracy is calculated from predicted coordinates. Coordinate accuracy alone cannot establish that the predictor ranks candidate outcomes correctly by familiarity. Measuring this ranking directly on shared contexts would test the prediction-to-memory interface more closely.

The biological interpretation concerns a proposed memory function rather than a complete cellular implementation. Spatial features are derived from supplied coordinates, and calcium recruits inhibition through a reduced pathway. The model does not establish how sensory representations would address astrocytic memory in a biological circuit. Differentiation through prediction heads, surrogate spike derivatives, and separate simulation clocks are additional implementation choices whose biological counterparts remain unspecified.

The evaluation is restricted to maps used for both tuning and training. New training seeds assess repeatability on those tasks, but do not test transfer to unfamiliar environments or new goals. Equal episode budgets also produce unequal numbers of observed transitions because successful episodes terminate earlier. Together with differences in learning rules and parameter counts between the recurrent and MLP controllers, this limits comparisons of learning efficiency. Evaluation on unseen maps, changed goals, and matched transition budgets would clarify how far the memory benefit extends beyond the conditions studied here.